\section{Migration Paths}
\label{sec:migration}

We now ask: starting from each configuration in
\S\ref{sec:configurations}, what transformations would be required to
reach a post-quantum-ready endpoint, and what blocks each
transformation? We categorize every transformation into one of four
blocker classes:

\begin{enumerate}
  \item[(a)] \textbf{Ready today.} The destination configuration is
    deployable; the transformation is engineering work against
    existing libraries and host functions.
  \item[(b)] \textbf{Blocked on host functions.} The destination is
    cryptographically defined and supported by mature libraries, but
    on the chosen anchor chain the verifier would run in pure
    bytecode at prohibitive gas cost. Unblocked by a chain-protocol
    upgrade adding the relevant precompile or host function.
  \item[(c)] \textbf{Blocked on tooling maturity.} The destination is
    cryptographically defined but lacks a production-grade
    implementation (no audited library, no compiler from the
    relation language to the proving system, no battle-tested setup
    procedure). Unblocked by engineering work outside the deploying
    team's scope.
  \item[(d)] \textbf{Blocked on cryptographic open problems.} The
    destination configuration would require a construction that does
    not yet exist in the literature, or one whose security is open.
    Unblocked only by research progress.
\end{enumerate}

Figure~\ref{fig:migration} renders the resulting migration DAG;
% Figure 2: Migration DAG for §6.
% Nodes: configurations from §5. Edges: feasible transformations.
% Edge labels: cost class (a/b/c/d per §6 categorization).
% PQ-ready endpoints highlighted.
%
% TODO: design after §5+§6 prose is drafted. Stub renders nothing.
\begin{figure}[t]
  \centering
  % \begin{tikzpicture}
  %   ...
  % \end{tikzpicture}
  \fbox{\parbox{0.9\columnwidth}{\centering\vspace{4em}TODO Figure 2:
    migration-path DAG\vspace{4em}}}
  \caption{Migration paths between configurations. Shaded nodes are
    PQ-ready endpoints. Edge labels denote blocker class. TODO redraw.}
  \label{fig:migration}
\end{figure}

each subsection discusses transformations from one starting point.

\subsection{From C1 (Stellar + BLS12-381 + Groth16)}
\label{sec:mig-c1}

Three migration paths to a PQ-ready endpoint:

\begin{description}
  \item[\textsc{C1 $\to$ C7} (in-place SNARK swap):] retain the
    Stellar anchor, swap Groth16 for Plonky3 / FRI. Class
    \textbf{(b)}: Plonky3 itself is mature enough for production
    use, but Stellar's host-function set --- BLS12-381 pairings
    (Protocol 22~\cite{stellar-cap-0059}), BN254 pairings (Protocol
    25~\cite{stellar-cap-0074}), and Poseidon / Poseidon2
    permutations over the existing scalar fields (also Protocol
    25~\cite{stellar-cap-0075}) --- does not include the FRI stack.
    CAP-0075 explicitly scopes out Goldilocks-field arithmetic, NTT,
    and FRI primitives; in pure WASM, FRI verification is one to two
    orders of magnitude more expensive than C1's pairing-precompile
    path. A successor CAP adding the missing FRI primitives would
    unblock this path; we discuss its technical shape in
    \S\ref{sec:open-host}.
  \item[\textsc{C1 $\to$ C6} (anchor change to StarkNet or similar):]
    abandon the Stellar anchor and re-deploy on a chain whose
    verification host is FRI-aware. Class \textbf{(a)}: the
    transformation is implementable today against existing tooling.
    Trade-offs: validator-set composition changes (StarkNet's prover
    is currently centralized; Stellar's $\sim$30-validator SCP is
    qualitatively different), fee market changes, and the team must
    re-audit against a different runtime.
  \item[\textsc{C1 $\to$ C1+lattice} (lattice-based pairings):]
    retain Stellar and Groth16-shaped relation but instantiate over a
    PQ-secure lattice-based pairing-replacement. Class \textbf{(d)}:
    no such construction exists.
\end{description}

\subsection{From C2 (Ethereum L1 + BN254 + Groth16)}
\label{sec:mig-c2}

\begin{description}
  \item[\textsc{C2 $\to$ STARK on L1} (in-place SNARK swap):] swap
    Groth16 for STARK on Ethereum L1. Class \textbf{(b)}:
    Ethereum L1 has BN254 precompiles (EIP-196/197) and is acquiring
    BLS12-381 precompiles via EIP-2537~\cite{eip-2537}, but no FRI
    host function is on the road-map; pure-bytecode verification is
    prohibitive at L1 gas prices.
  \item[\textsc{C2 $\to$ C6} (anchor change to StarkNet):] re-deploy
    on a FRI-aware L2. Class \textbf{(a)}. Trade-offs: introduces the
    L2 sequencer trust layer; the L1-anchor-only fallback present in
    rollup architectures partially mitigates this.
  \item[\textsc{C2 $\to$ Stellar} (anchor change with curve
    preservation):] re-deploy a Tornado-style construction on
    Stellar without curve change. Class \textbf{(a)} as of Stellar
    Protocol 25, which added native BN254 host
    functions~\cite{stellar-cap-0074} (point addition, scalar
    multiplication, pairing check). A C2 deployer can keep the
    BN254 curve, the BN254-instantiated Groth16 verifying key, the
    Poseidon arithmetization (CAP-0075~\cite{stellar-cap-0075}),
    and the per-circuit MPC ceremony output, and only swap the
    anchor chain and the contract host. Does not improve PQ
    posture; included as a non-PQ migration that changes only the
    trust posture and the fee market.
\end{description}

\subsection{From C3 (Aztec + UltraPlonk)}
\label{sec:mig-c3}

\begin{description}
  \item[\textsc{C3 $\to$ STARK-Aztec} (in-place SNARK family swap):]
    re-implement the Aztec rollup machinery on a hash-based proof
    system, retaining the verifier-as-rollup posture. Class
    \textbf{(c)}: UltraPlonk-on-BN254 is deeply load-bearing in the
    Aztec architecture (the Noir compiler, the proving server, the
    sequencer's batched-proof aggregator); a STARK-based replacement
    is engineering scale outside the registry deployer's scope.
  \item[\textsc{C3 $\to$ C6} (anchor change):] migrate the registry's
    relations to a STARK-based L2 (StarkNet). Class \textbf{(a)} for
    the relation re-implementation, but the application loses
    Aztec's privacy-as-execution-model semantics. The PQ posture
    improves; the privacy posture changes shape.
\end{description}

\subsection{From C4 (Mina + Pickles + Pasta)}
\label{sec:mig-c4}

\begin{description}
  \item[\textsc{C4 $\to$ FRI-Pickles} (in-place commitment-scheme
    swap):] replace Pasta-cycle commitments with a hash-based
    commitment scheme. Class \textbf{(c)}: Mina's road-map
    discusses this (Kimchi has a PQ trajectory) but no
    production-grade implementation exists. The recursion structure
    of Pickles is an active research area for hash-based
    re-instantiation.
  \item[\textsc{C4 $\to$ C6} (anchor change):] re-deploy on
    StarkNet. Class \textbf{(a)} for the registry, but the
    application loses the protocol-level recursion that makes Mina
    distinctive.
\end{description}

\subsection{From C5 (L2 + Plonky2/3 + KZG over BN254)}
\label{sec:mig-c5}

C5 is the most interesting starting point because the inner prover is
already FRI-based; only the outer KZG wrap is post-quantum-broken.

\begin{description}
  \item[\textsc{C5 $\to$ STARK-wrapped C5} (outer-wrap swap):]
    replace the KZG wrap with a STARK wrap and submit the wrapped
    proof to the L1 verifier. Class \textbf{(b)}: Ethereum L1 lacks
    a FRI host function; the wrapped proof verifies in pure bytecode
    at uneconomical gas. On an L1 with FRI host functions, this is
    Class (a).
  \item[\textsc{C5 $\to$ C6} (drop the L1 anchor):] settle directly
    on a FRI-aware L2 such as StarkNet, eliminating the need for the
    KZG wrap entirely. Class \textbf{(a)}. The L1-anchor-only
    fallback property is lost; configurations that depend on it must
    decline this migration.
\end{description}

\subsection{From C6 (STARK over Goldilocks)}
\label{sec:mig-c6}

C6 is already PQ-ready. Migrations from C6 are about
\emph{portability} rather than PQ posture: which chains can host the
configuration?

\begin{description}
  \item[\textsc{C6 portability to Stellar}:] equivalent to the
    \textsc{C1 $\to$ C7} path. Class \textbf{(b)}.
  \item[\textsc{C6 portability to Ethereum L1}:] equivalent to the
    \textsc{C2 $\to$ STARK-on-L1} path. Class \textbf{(b)}.
  \item[\textsc{C6 portability to FRI-aware L2}:] Class \textbf{(a)};
    StarkNet, Polygon zkEVM with FRI inner layer, and similar chains
    accept the configuration directly.
\end{description}

\subsection{From C7 (hypothetical Stellar + Plonky3 + FRI)}
\label{sec:mig-c7}

C7 is the asymmetric case: it is post-quantum-ready and would inhabit
the host-function-cheap regime were it deployable, but it is not. The
sole migration is the host-function unblock.

\begin{description}
  \item[\textsc{Stellar protocol upgrade adding FRI host functions}:]
    Class \textbf{(b)}. Protocol 25 has already moved part of the
    way: CAP-0075~\cite{stellar-cap-0075} added Poseidon and
    Poseidon2 permutation primitives over the existing scalar fields
    (BLS12-381~$\mathbb{F}_r$, BN254~$\mathbb{F}_r$). What remains
    missing for a deployable FRI verifier is the orthogonal subset
    that CAP-0075 explicitly declined to introduce: Goldilocks-field
    arithmetic, NTT, and FRI-specific operations over a small field.
    The CAP's rationale documents this scope decision; it is
    deliberate rather than oversight. A successor CAP would need to
    introduce a new field arithmetic to the host-function ABI, which
    is a larger commitment than CAP-0075 was willing to make. We
    estimate the engineering scale of the unblock at one to two
    chain-protocol-upgrade cycles plus validator activation; we do
    not estimate the political feasibility, which is the
    load-bearing variable in practice.
\end{description}

\subsection{Migration matrix}
\label{sec:migration-matrix}

\begin{table}[t]
  \centering
  \caption{Migration blocker class for each (start, end) pair where
    the end is a PQ-ready configuration. ``--'' indicates the
    transformation is undefined or trivially the identity.
    Table~\ref{tab:migration} populates the off-diagonal cells of
    \S\ref{sec:configurations}.}
  \label{tab:migration}
  \begin{tabular}{lccc}
    \toprule
    Start & $\to$ C6 & $\to$ C7 & lattice path \\
    \midrule
    C1 & (a) & (b) & (d) \\
    C2 & (a) & ---  & (d) \\
    C3 & (a)$^\ast$ & --- & (d) \\
    C4 & (a)$^\ast$ & (c)$^\dagger$ & (d) \\
    C5 & (a) & (b) & (d) \\
    C6 & ---  & (b) & (d) \\
    C7 & (b) & ---  & (d) \\
    \bottomrule
  \end{tabular}
  \\[0.5em]
  \footnotesize
  $^\ast$ The migration is technically Class (a) but loses
    application-relevant semantics (verifier-as-rollup for C3,
    protocol-level recursion for C4). \\
  $^\dagger$ FRI-Pickles is the C4 in-place migration; included
    for completeness in the C7-shaped destination column.
\end{table}

The matrix carries one observation that drives \S\ref{sec:open}: no
column is uniformly Class (a). Either the migration is operational
(anchor change), losing some property; or it is blocked on host
functions; or it is blocked on tooling; or it is open. The
deployment recommendations of \S\ref{sec:recommendations} navigate
this landscape by starting position and priority weighting rather
than by attempting to identify a globally dominant migration.
